\chapter{Transurethral Resection of the Prostate (TURP)}

\section*{Introduction \& Clinical Indications}
Transurethral Resection of the Prostate (TURP) is a minimally invasive surgical procedure designed to alleviate urinary obstruction caused by benign prostatic hyperplasia (BPH). This condition, characterized by the non-cancerous enlargement of the prostate gland, can lead to significant lower urinary tract symptoms (LUTS) such as urinary frequency, urgency, and nocturia. TURP involves the endoscopic removal of obstructing prostatic tissue through the urethra, utilizing a specialized instrument called a resectoscope. The primary goal of this procedure is to restore normal urinary flow by excising the hyperplastic tissue that compresses the prostatic urethra. TURP remains a cornerstone in the surgical management of BPH, offering symptomatic relief and improved quality of life for many patients.

\begin{itemize}
  \item Endoscopic Prostatectomy
  \item Resection of the Prostate
  \item Transurethral Prostatectomy
  \item Prostate Scrape
  \item BPH Surgery
  \item Conventional TURP
\end{itemize}

The surgical principle of TURP is based on the excision of obstructive prostatic tissue to relieve bladder outlet obstruction. The procedure employs a resectoscope equipped with an electrosurgical loop that cuts and coagulates tissue simultaneously. The mechanism of correction involves the removal of adenomatous tissue that compresses the prostatic urethra, thereby restoring unobstructed urinary flow. 

The technical goals of TURP include:

\begin{itemize}
  \item Achieving a wide prostatic fossa to facilitate normal urinary passage.
  \item Minimizing blood loss through effective hemostasis during resection.
  \item Ensuring complete removal of obstructive tissue to alleviate symptoms of LUTS.
\end{itemize}

By addressing the mechanical obstruction, TURP significantly improves urinary flow rates and reduces post-void residual urine volumes, enhancing the patient's quality of life.

The evolution of TURP can be traced back to the early 20th century, with significant milestones marking its development. In 1909, Hugh Hampton Young introduced the punch resectoscope, which laid the groundwork for endoscopic prostatectomy. However, early techniques were limited by inadequate visualization and hemostasis.

A pivotal advancement occurred in 1926 when Maximilian Stern developed the first resectoscope with high-frequency cutting current, allowing for more effective tissue removal. Joseph F. McCarthy further refined the technique in 1931 by introducing a continuous irrigation system, enhancing safety and efficacy.

Over the decades, TURP became the gold standard for surgical treatment of BPH, celebrated for its effectiveness and minimally invasive nature compared to open prostatectomy. Despite the emergence of newer techniques such as laser prostatectomy and minimally invasive surgical therapies, TURP remains a widely performed procedure, particularly for patients with moderate-sized prostates.

TURP is indicated for patients experiencing significant lower urinary tract symptoms (LUTS) due to benign prostatic hyperplasia (BPH) when conservative management has failed. Specific clinical indications include:

\begin{itemize}
  \item Moderate-to-severe LUTS that impair quality of life and are refractory to medical therapy.
  \item Recurrent acute urinary retention that cannot be managed by catheterization alone.
  \item Recurrent urinary tract infections secondary to large post-void residual urine volumes.
  \item Formation of bladder stones due to urinary stasis from prostatic obstruction.
  \item Gross hematuria originating from prostatic congestion unresponsive to conservative measures.
  \item Renal insufficiency or hydronephrosis caused by chronic bladder outlet obstruction.
  \item Significant post-void residual urine volume indicating inefficient bladder emptying.
\end{itemize}

These indications highlight the importance of TURP in managing complications arising from BPH and improving patient outcomes.

TURP serves as both a primary and secondary surgical intervention for benign prostatic hyperplasia. It is considered a primary option for patients with moderate to severe symptoms or complications of BPH who have not responded to medical therapy. In this context, TURP is often the first surgical intervention considered.

However, with the advent of various minimally invasive surgical therapies (MISTs) such as prostatic urethral lift (UroLift) and water vapor thermal therapy (Rezum), TURP may also be viewed as a secondary option in certain scenarios. It is particularly relevant when less invasive options are unsuitable due to prostate size, anatomy, or patient comorbidities. Despite the rise of alternatives, TURP remains a highly effective and durable procedure, especially when a more definitive and complete resection is desired.

\section*{Relevant Surgical Anatomy}
The prostate gland is a walnut-sized organ located inferior to the bladder and surrounds the prostatic urethra. It is divided into several zones, with the transition zone being the primary site for hyperplastic growth in BPH. The arterial supply to the prostate is primarily from the inferior vesical and middle rectal arteries, while venous drainage occurs through the prostatic venous plexus, which drains into the internal iliac veins. 

Key anatomical landmarks relevant to TURP include:

\begin{itemize}
  \item **Bladder Neck:** The junction where the bladder meets the urethra, often the site of obstruction in BPH.
  \item **Verumontanum:** A prominent anatomical landmark located in the prostatic urethra, serving as a reference point during resection.
  \item **External Urethral Sphincter:** Important for urinary continence, located just distal to the prostate.
  \item **Seminal Vesicles:** Located posterior to the prostate, these structures should be preserved during TURP to maintain fertility.
\end{itemize}

Understanding these anatomical relationships is crucial for the safe and effective performance of TURP, minimizing the risk of complications such as bleeding or injury to surrounding structures.

\section*{Preoperative Workup \& Preparation}
Thorough preoperative preparation is essential for a successful TURP. The preoperative workup includes:

\begin{itemize}
  \item **Medical History and Physical Exam:** Comprehensive assessment of overall health and identification of comorbidities.
  \item **Laboratory Tests:** Complete blood count (CBC), basic metabolic panel, coagulation studies, prostate-specific antigen (PSA), and urinalysis with culture to rule out urinary tract infections.
  \item **Imaging Studies:** Renal ultrasound to assess for hydronephrosis and measurement of post-void residual (PVR) volume.
  \item **Cardiac and Anesthesia Clearance:** Clearance from specialists for patients with significant cardiac or pulmonary conditions.
  \item **Medication Adjustments:** Discontinuation of anticoagulant or antiplatelet medications several days prior to surgery.
  \item **NPO Status:** Patients must fast for at least 6-8 hours before surgery.
  \item **Antibiotic Prophylaxis:** Administration of prophylactic antibiotics just before the procedure.
  \item **Informed Consent:** Detailed discussion of the procedure, risks, benefits, and alternatives with the patient.
\end{itemize}

This comprehensive preparation helps minimize complications and optimize surgical outcomes.

While TURP is a common and effective procedure, certain conditions may contraindicate its performance or necessitate caution.

**Absolute Contraindications:**
\begin{itemize}
  \item Active, untreated urinary tract infection (UTI) or prostatitis.
  \item Severe, uncorrectable coagulopathy or bleeding disorder.
  \item Known prostate cancer requiring radical prostatectomy.
  \item Severe bladder dysfunction where obstruction is not the primary cause of voiding difficulty.
  \item Urethral stricture preventing passage of the resectoscope.
\end{itemize}

**Relative Contraindications and Precautions:**
\begin{itemize}
  \item Very large prostate gland (>80-100 grams).
  \item Severe cardiovascular or pulmonary disease.
  \item Prior pelvic radiation therapy.
  \item Neurogenic bladder dysfunction.
  \item Uncontrolled diabetes mellitus.
  \item Patients desiring to preserve ejaculatory function.
\end{itemize}

Recognizing these contraindications is crucial for patient safety and optimal surgical outcomes.

\section*{Surgical Technique \& Procedure}
The TURP procedure is typically performed under general or spinal anesthesia. The patient is positioned in the dorsal lithotomy position, and the perineum and external genitalia are prepped and draped in a sterile manner. The resectoscope is then carefully inserted through the urethra into the bladder.

Key steps in the TURP procedure include:

\begin{itemize}
  \item **Visualization and Landmarks:** Identification of anatomical landmarks such as the bladder neck, verumontanum, and external urethral sphincter.
  \item **Resection of Prostatic Lobes:** The surgeon uses the electrosurgical wire loop to shave off small chips of prostatic tissue, typically starting with the median lobe, followed by the lateral lobes.
  \item **Hemostasis:** Continuous coagulation of bleeding vessels is performed to maintain a clear surgical field and minimize blood loss.
  \item **Chip Evacuation:** Resection chips are evacuated from the bladder using continuous irrigation with a non-electrolytic solution, usually glycine.
  \item **Completion and Inspection:** A final inspection is conducted to ensure complete hemostasis and an unobstructed channel.
  \item **Catheter Placement:** A three-way Foley catheter is inserted for continuous bladder irrigation to prevent clot formation and facilitate drainage.
\end{itemize}

The procedure typically lasts 60-90 minutes, depending on the size of the prostate and the extent of resection required.

\section*{Complications \& Risks}
As with any surgical procedure, TURP carries potential risks and complications, categorized as follows:

\begin{itemize}
  \item **Anesthesia-related Risks:**
    \begin{itemize}
      \item Adverse reactions to anesthetic medications.
      \item Respiratory or cardiac complications.
      \item Nerve damage or spinal headache (with spinal anesthesia).
    \end{itemize}
  
  \item **General Surgical Risks:**
    \begin{itemize}
      \item **Bleeding:** Significant intraoperative or postoperative hemorrhage.
      \item **Infection:** Urinary tract infection, epididymitis, prostatitis, or sepsis.
      \item **Pain:** Postoperative discomfort at the surgical site.
    \end{itemize}
  
  \item **Procedure-Specific Risks:**
    \begin{itemize}
      \item **Retrograde Ejaculation:** Occurs in 70-90\% of patients, often permanent.
      \item **TURP Syndrome:** Rare but serious complication from excessive absorption of irrigation fluid.
      \item **Urethral Stricture:** Scarring and narrowing of the urethra postoperatively.
      \item **Bladder Neck Contracture:** Scarring at the bladder-prostate junction.
      \item **Urinary Incontinence:** Temporary or persistent leakage of urine.
      \item **Erectile Dysfunction:** New or worsened erectile dysfunction in a small percentage of patients.
      \item **Perforation of the Prostatic Capsule or Bladder:** Accidental injury requiring repair.
      \item **Retained Prostatic Chips:** Incomplete evacuation leading to recurrent obstruction.
    \end{itemize}
\end{itemize}

Awareness of these complications allows for better patient counseling and management strategies.

\section*{Postoperative Care \& Recovery}
Postoperative recovery following TURP typically involves several phases. Immediately after surgery, patients are monitored in the post-anesthesia care unit (PACU) for vital signs, pain levels, and urine output. A three-way Foley catheter is placed for continuous bladder irrigation, which helps prevent clot formation and maintains urine drainage.

In the first 24-48 hours, the continuous bladder irrigation is adjusted to maintain clear urine output. The catheter is usually removed within 1-3 days, depending on the patient's recovery and urine clarity. Patients may experience dysuria, frequency, and urgency as the prostatic fossa heals. 

During the first 1-2 weeks, patients are encouraged to gradually resume normal activities while avoiding heavy lifting and strenuous exercise. Long-term recovery may take several weeks to months, with gradual improvement in urinary symptoms as healing progresses.

During TURP, the surgeon documents the extent of resection and any intraoperative findings. The pathology report on the resected tissue typically confirms the diagnosis of benign prostatic hyperplasia and rules out malignancy. This report is crucial for guiding further management and monitoring.

A successful postoperative course following TURP is characterized by significant improvement in lower urinary tract symptoms (LUTS). Patients typically experience a gradual resolution of dysuria and other irritative symptoms within weeks. The expected timeline for symptom improvement includes increased urinary flow rates and reduced post-void residual volumes. Most patients report a marked enhancement in quality of life and satisfaction with the procedure.

Postoperative care is essential for optimal recovery following TURP. Key components include:

\begin{itemize}
  \item **Postoperative Visits:** Follow-up appointments are scheduled 2-4 weeks after surgery to assess healing and discuss pathology results.
  \item **Activity Restrictions:** Patients should avoid strenuous activities and heavy lifting for 4-6 weeks.
  \item **Dietary Advice:** A high-fiber diet and adequate fluid intake are recommended to prevent constipation.
  \item **Medication Review:** Adjustments to medications for bladder symptoms may be necessary.
  \item **Monitoring for Complications:** Patients should be vigilant for warning signs such as persistent bleeding, fever, or inability to urinate.
  \item **Uroflowmetry and PVR Testing:** These tests may be repeated to assess improvement in urinary function.
  \item **PSA Monitoring:** Annual PSA testing may continue as part of routine prostate cancer screening.
\end{itemize}

Long-term follow-up is crucial to ensure sustained symptom relief and address any potential late complications.

\section*{Outcomes \& Prognosis}
Benign prostatic hyperplasia (BPH) is a prevalent condition, particularly among older men. Microscopic evidence of BPH is found in approximately 50\% of men in their 50s, 70\% in their 60s, and 90\% in their 80s. Clinically significant BPH requiring intervention affects about 25-30\% of men by age 60 and up to 50\% by age 80. TURP remains one of the most common surgical procedures for BPH, with hundreds of thousands performed annually worldwide. The typical age range for TURP patients is 60-80 years. Mortality rates associated with TURP are low, generally less than 0.2\%, while morbidity rates can include complications such as retrograde ejaculation (70-90\%) and bleeding requiring transfusion (5-10\%).

The outcomes of TURP are generally favorable, with success rates for relieving lower urinary tract symptoms (LUTS) ranging from 80-90\%. Patients typically experience significant improvements in peak urinary flow rates and reductions in post-void residual volumes. Long-term prognosis is good, with most patients enjoying sustained relief for many years. However, recurrence of BPH symptoms requiring reoperation can occur in approximately 10-20\% of patients within 10 years. Factors influencing outcomes include initial prostate size, severity of preoperative symptoms, and presence of detrusor dysfunction. While TURP effectively addresses mechanical obstruction, some irritative symptoms may persist in a subset of patients.

\section*{References}
\begin{enumerate}
  \item MedlinePlus. Transurethral Resection of the Prostate (TURP). U.S. National Library of Medicine; 2024.
  \item Wein AJ, Kavoussi LR, Partin AW, Peters CA, editors. Campbell-Walsh-Wein Urology. 12th ed. Elsevier; 2021.
  \item American Urological Association (AUA). AUA Guideline on the Management of Benign Prostatic Hyperplasia (BPH). 2021.
  \item Roehrborn CG. BPH: Epidemiology and natural history. Urol Clin North Am. 2004;31(3):385-399.
\end{enumerate}

\clearpage
